\section{Model, state spaces, and conventions}

\subsection{Reduced momentum and physical wavevector}

The Brillouin coordinate used by every installed numerical owner is the dimensionless
reduced crystal momentum
\begin{equation}
 q\in\Brillouin=[-\tfrac12,\tfrac12),\qquad q\sim q+1.
\end{equation}
If $a_0$ is a physical lattice period, $G=2\pi/a_0$ is the reciprocal-lattice period
and $k$ is the inverse-length wavevector, then
\begin{equation}
 q=\frac{k}{G}=\frac{k a_0}{2\pi},\qquad k=qG.
 \label{eq:reduced-physical-relation}
\end{equation}
Accordingly, all band functions in this report are written as $\lambda_n(q)$.
The symbol $k$ is reserved for a physical wavevector and is never used as a finite-mesh
index.  Integer mesh indices are denoted $j$, amplitude indices by $r$, and cell
displacements by $R$.

For even $N$, the centered half-open construction mesh is
\begin{equation}
 \mathcal Q_N=\left\{q_j=-\frac12+\frac{j}{N}:j=0,\ldots,N-1\right\}.
 \label{eq:centered-mesh}
\end{equation}
The half-open convention avoids duplicating the periodically equivalent endpoints.
All interval statements involving $q$ are understood modulo one.

\subsection{Ambient and retained state spaces}

The ambient cell space is $\Complex^4$ in a fixed ordered four-orbital basis.  The
retained fiber at momentum $q$ is a two-dimensional spectral subspace
$\mathcal S(q)\subset\Complex^4$.  A frame
$F(q)\in\Complex^{4\times2}$ is an isometric coordinate map from the retained
coordinate space $\Complex^2$ into $\mathcal S(q)$:
\begin{equation}
 F(q)^\dagger F(q)=I_2,
 \qquad
 P(q)=F(q)F(q)^\dagger.
 \label{eq:frame-projector}
\end{equation}
The physical subspace is represented by the projector $P(q)$ and is unchanged by a
right unitary action $F(q)\mapsto F(q)W(q)$.  The represented $2\times2$ Hamiltonian,
however, changes by conjugation.  This distinction between subspace and representation
is the central conceptual layer of the study.

\subsection{Finite-range four-orbital parent}

Two continuous maps $f_0,f_1:X\to Y$ are \emph{homotopic} when there exists a
continuous map
\begin{equation}
 \mathcal H:X\times[0,1]\longrightarrow Y
 \quad\text{such that}\quad
 \mathcal H(x,0)=f_0(x),\qquad \mathcal H(x,1)=f_1(x).
 \label{eq:homotopy-definition}
\end{equation}
The map $\mathcal H$ is called a \emph{homotopy} between $f_0$ and $f_1$.  In this
study, $X=\Brillouin$ is the period-one Brillouin circle, $Y=\operatorname{Herm}(4)$ is
the space of $4\times4$ Hermitian matrices, and the interpolation parameter is
$\rho\in[0,1]$.  Consequently, if $(q,\rho)\mapsto K_\rho(q)$ is continuous, then the
endpoint parent fields $K_0$ and $K_1$ are homotopic as Hermitian matrix-valued maps.
This use of \emph{homotopy} means a continuous parent-model interpolation only.  It does
not by itself establish a gap-preserving homotopy or topological equivalence: that would
require the relevant spectral gap to remain open for every continuous pair
$(q,\rho)$, whereas the calculation applies its external-gap gate only on finite sampled
inventories.

Let $z(q)=e^{2\pi iq}$.  The installed parent-hopping homotopy is the continuous map
\begin{equation}
 K_\rho(q)=
 \begin{pmatrix}
 L_0+\rho[z(q)L_1+\overline{z(q)}L_1^\dagger]&\rho z(q)C\\
 \rho\overline{z(q)}C^\dagger&
 U_0+\rho[z(q)U_1+\overline{z(q)}U_1^\dagger]
 \end{pmatrix},
 \qquad 0\le\rho\le1,
 \label{eq:installed-parent}
\end{equation}
where
\begin{align}
 L_0&=\begin{pmatrix}-1.20&0.06\\0.06&-0.35\end{pmatrix},&
 L_1&=\begin{pmatrix}0.15&0\\0&-0.10\end{pmatrix},\\
 U_0&=\begin{pmatrix}1.00&0.04\\0.04&1.75\end{pmatrix},&
 U_1&=\begin{pmatrix}0.05&0\\0&-0.04\end{pmatrix},\\
 C&=\begin{pmatrix}0.12&0\\0&0.10\end{pmatrix}.
 \label{eq:installed-blocks}
\end{align}
All numerical entries use one arbitrary positive synthetic parent-energy unit.  There
is no conversion to electronvolts and no material parametrization.  Equation
\eqref{eq:installed-parent} is exactly Hermitian because its lower-left block is the
adjoint of its upper-right block and the diagonal blocks include conjugate hopping
terms.

Equivalently,
\begin{equation}
 K_\rho(q)=\sum_{R=-1}^{1}e^{2\pi iqR}T_R(\rho),
 \qquad T_{-R}(\rho)=T_R(\rho)^\dagger.
 \label{eq:finite-parent-series}
\end{equation}
At $\rho=0$ the ambient parent is onsite and independent of $q$.  At $\rho=1$ it is the
full installed finite-range parent.  Intermediate $\rho$ values vary only the non-onsite
parent hopping; the onsite blocks, state spaces, gates, tolerances, and gauge-amplitude
inventory remain unchanged.

\subsection{Spectral retention and external-gap gate}

Let
\begin{equation}
 \lambda_1(q)\le\lambda_2(q)\le\lambda_3(q)\le\lambda_4(q)
\end{equation}
be the ordered eigenvalues of $K_\rho(q)$.  The raw retained frame $V_j$ consists of the
first two normalized eigenvectors returned at $q_j$.  The sampled external isolation
gap is
\begin{equation}
 \Delta_{\mathrm{ext},N}(\rho)
 =\min_{q_j\in\mathcal Q_N}\left[\lambda_3(q_j)-\lambda_2(q_j)\right].
 \label{eq:external-gap}
\end{equation}
Retention is allowed to proceed only if
$\Delta_{\mathrm{ext},N}(\rho)\ge\delta_{\mathrm{ext}}$.  Failure of this gate is a
scientific qualification stop, not a numerical transport failure and not a failed
comparison threshold.

The raw eigenvectors are not presumed to define a smooth periodic gauge.  Arbitrary
eigenvector phases and rotations inside the rank-two eigenspace must be resolved before
Fourier locality is meaningful.  Section~3.2 specifies the closed-loop polar transport
used for that purpose.

\subsection{Represented retained Hamiltonian}

Once a periodic transported frame $F_j=F(q_j)$ is available, the retained Hamiltonian is
\begin{equation}
 H_j=F_j^\dagger K_\rho(q_j)F_j\in\Complex^{2\times2}.
 \label{eq:retained-hamiltonian}
\end{equation}
Its scalar and traceless parts are
\begin{equation}
 m_j=\frac12\tr H_j,
 \qquad X_j=H_j-m_jI_2,
 \qquad \tr X_j=0.
 \label{eq:trace-traceless}
\end{equation}
The installed splitting scale is the discrete normalized traceless power
\begin{equation}
 \kappa_N=\frac{1}{2N}\sum_{j=0}^{N-1}\norm{X_j}_F^2>0.
 \label{eq:splitting-scale}
\end{equation}
Division by two is division by retained rank.  Thus every squared Frobenius quantity in
the main comparison has the dimensions of energy squared per retained state.

\subsection{Gauge attack and exact pointwise band equivalence}

The declared smooth gauge angle and real rank-two rotation are
\begin{equation}
 \theta_a(q)=a\sin(2\pi q),
 \qquad
 \mathcal R(\theta)=
 \begin{pmatrix}\cos\theta&-\sin\theta\\
                 \sin\theta& \cos\theta\end{pmatrix},
 \qquad 0\le a\le\frac{\pi}{4}.
 \label{eq:gauge-rotation}
\end{equation}
The comparison field is
\begin{equation}
 A_a(q)=\mathcal R(\theta_a(q))^\dagger
        H(q)\mathcal R(\theta_a(q)).
 \label{eq:comparison-field}
\end{equation}
Because $\mathcal R$ is orthogonal and hence unitary,
\begin{equation}
 \lambda_n[A_a(q)]=\lambda_n[H(q)],\qquad n=1,2,
 \label{eq:exact-band-equivalence}
\end{equation}
mathematically in exact arithmetic.  The finite calculation nevertheless retains the
maximum observed sorted-band defect as an implementation residual.

The mean energy $m(q)$ is invariant under conjugation.  Defining
$Y_a(q)=A_a(q)-m(q)I_2$, one has
$Y_a=\mathcal R^\dagger X\mathcal R$ and
$\norm{Y_a(q)}_F=\norm{X(q)}_F$ pointwise.  The transformation therefore redistributes
matrix orientation but does not alter the local eigenvalue splitting.

\subsection{Flat analytic control versus embedded dispersive control}

The flat analytic control is a separate retained-space model
\begin{equation}
 H_{\mathrm{flat}}(q)=\bar E I_2+D,
 \qquad D=D^T\in\Reals^{2\times2},
 \quad\tr D=0,
 \quad D\ne0.
 \label{eq:flat-control}
\end{equation}
It has no parent dispersion and no non-onsite reference hopping.  Under the sinusoidal
gauge attack, its continuum restricted correlation has a Bessel-function expression
derived in Section~3.4.  The embedded control of
Eqs.~\eqref{eq:installed-parent}--\eqref{eq:retained-hamiltonian} is different: its
retained field is dispersive, its frame comes from an ambient spectral subspace, and its
restricted correlation has no assumed $J_0$ identity.

These layers must not be conflated.  The flat calculation is an analytic control; the
$\rho$ family is a parent-model robustness channel; and mesh refinement is a numerical
discretization channel.
